\section{Per-Chain Analysis}


\subsection{Bitcoin}

\subsubsection{Key Generation and Address Derivation}

Bitcoin generates an ECDSA key pair $(sk, pk)$ on secp256k1. The
address is derived through two hash layers:

\begin{enumerate}[nosep]
  \item $h_1 = \text{SHA-256}(pk)$
  \item $h_{160} = \text{RIPEMD-160}(h_1)$
  \item $\text{address} = \text{Base58Check}(\text{version} \| h_{160}
        \| \text{checksum})$
\end{enumerate}

\subsubsection{Public Key Exposure}

The public key $pk$ is \emph{not} revealed when receiving Bitcoin. The
address is a hash of $pk$, and hash preimage resistance protects $pk$
even if the address is known. The public key is revealed only when
\emph{spending}---the transaction includes $pk$ for signature
verification.

\subsubsection{Quantum Vulnerability Window}

The exposure window is the time between transaction broadcast and
block confirmation ($\sim$10 minutes for one confirmation, $\sim$60
minutes for six). An adversary must derive $sk$ from $pk$ within this
window to redirect funds.

\begin{remark}
Bitcoin's recommended practice of single-use addresses provides
additional protection: once an address's public key is revealed by
spending, no further funds should be sent to it. Unspent outputs at
never-used addresses have zero public key exposure.
\end{remark}

\subsection{Ethereum}

\subsubsection{Key Generation and Address Derivation}

Ethereum generates an ECDSA key pair on secp256k1. The address is
the last 20 bytes of the Keccak-256 hash of the public key:

\[
\text{address} = \text{Keccak-256}(pk)[12{:}32]
\]

\subsubsection{Public Key Exposure}

The public key is \textbf{effectively always exposed}. It can be
recovered from any ECDSA signature via ecrecover:

\[
pk = \text{ecrecover}(\text{msgHash}, v, r, s)
\]

Once an account has signed any transaction, the public key is
permanently derivable from the blockchain history.

\subsubsection{Quantum Vulnerability Window}

The exposure window is \textbf{permanent} for any account that has
ever sent a transaction. An adversary does not need to race against
block confirmation---they can extract $pk$ from historical transactions
at leisure and compute $sk$ offline.

\begin{remark}
Ethereum's account model exacerbates the risk. Unlike Bitcoin's UTXO
model where each output can use a fresh key, Ethereum accounts are
long-lived and typically reuse the same key for all transactions.
\end{remark}

\subsection{Solana}

\subsubsection{Key Generation and Address Derivation}

Solana uses Ed25519 key pairs. Unlike Bitcoin (SHA-256 + RIPEMD-160)
or Ethereum (Keccak-256), Solana applies \emph{no hash layer at all}:
the address is the raw public key, base58-encoded.

\[
\text{address} = \text{base58}(pk)
\]

\subsubsection{Public Key Exposure}

The public key is \emph{always} visible to anyone who knows the
address---there is no hash-based protection, even before first spend.
This is strictly weaker than Bitcoin's UTXO model (pk hidden until
spend) and equivalent in outcome to Ethereum's permanent exposure.

\subsubsection{Quantum Vulnerability Window}

Permanent from the moment an address is published. A quantum adversary
with Shor's algorithm can recover the secret key from the address
directly, without needing a prior transaction. Solana has the
\emph{worst} public-key exposure profile of the four chains analysed
here.

\subsection{Cardano}

\subsubsection{Key Generation and Address Derivation}

Cardano uses separate spending and staking key pairs. Addresses are
derived via Blake2b-224:

\begin{align*}
\text{PaymentKeyHash} &= \text{Blake2b-224}(\text{PublicSpendingKey}) \\
\text{StakeKeyHash} &= \text{Blake2b-224}(\text{PublicStakingKey}) \\
\text{BaseAddress} &= \text{Bech32}(\text{NetworkID} \|
  \text{PaymentKeyHash} \| \text{StakeKeyHash})
\end{align*}

\subsubsection{Public Key Exposure}

The spending public key is revealed when signing a transaction.
Cardano's extended UTXO model allows address reuse, but best practice
is to derive fresh addresses for each transaction.

\subsubsection{Quantum Vulnerability Window}

Similar to Bitcoin: the public key is exposed at signing time. However,
Cardano's faster block times ($\sim$20 seconds) reduce the exposure
window compared to Bitcoin's 10 minutes.

%% ============================================================
